% SEGMENT OLP-0626-S01
% Part: history
% Chapter: biographies
% Section: emmy-noether

\documentclass[../../../include/open-logic-section]{subfiles}

\begin{document}

\olfileid{his}{bio}{noe}

\olsection{எம்மி நேதர்}

\olphoto{noether-emmy}{எம்மி நேதர்}

எம்மி நேதர் (\textsc{ner}-ter) 1882 மார்ச் 23 அன்று ஜெர்மனியின்
எர்லாங்கனில், கல்வியில் ஈடுபட்டிருந்த மேல் நடுத்தரக் குடும்பத்தில்
பிறந்தார். ``நவீன இயற்கணிதத்தின் தாய்'' என்று போற்றப்படும் நேதர்,
பெண்களின் கல்விக்கு இருந்த பெரிய தடைகளையும் மீறிக் கணிதத்திற்கும்
இயற்பியலுக்கும் புரட்சிகரமான பங்களிப்புகளைச் செய்தார். அக்கால
ஜெர்மனியில் சிறுமிகளுக்குக் கலைகளில் கல்வி கற்பிப்பதே வழக்கமாக இருந்தது;
கல்லூரிக்கான ஆயத்தப் பள்ளிகளில் சேர அவர்களுக்கு அனுமதி இல்லை. எனினும்,
க\"{o}ட்டிங்கன், எர்லாங்கன் பல்கலைக்கழகங்களில் வகுப்புகளில்
பதிவுசெய்யாமலே பங்கேற்ற பின்னர், 1904-இல் எர்லாங்கன் பெண்களை மாணவர்களாக
அனுமதிக்கத் தனது கொள்கையை மாற்றியபோது அங்கே சேர முடிந்தது. அவருடைய
தந்தை அங்கே கணிதப் பேராசிரியராக இருந்தார். 1907-இல் கணிதத்தில்
முனைவர் பட்டம் பெற்றார்.

% SEGMENT OLP-0626-S02
தகுதிகள் இருந்தபோதிலும் நேதர் தமது பணிவாழ்வில் மிகுந்த எதிர்ப்பைச்
சந்தித்தார். 1908 முதல் 1915 வரை எர்லாங்கனில் ஊதியமின்றிக் கற்பித்தார்.
அக்காலத்தின் தலைசிறந்த கணிதவியலாளர்களில் ஒருவரான டேவிட் ஹில்பர்ட்
அவருடைய திறமையைக் கவனித்து க\"{o}ட்டிங்கனுக்கு அழைத்தார். ஆனால் பெண்கள்
பேராசிரியர் பதவி பெறத் தடை இருந்ததால், ஹில்பர்ட்டின் பெயரில் மட்டுமே
அவரால் உரையாற்ற முடிந்தது; அதற்கும் ஊதியம் இல்லை. இக்காலத்தில்,
இன்றும் கோட்பாட்டு இயற்பியலில் பயன்படும் நேதரின் தேற்றம் எனப்படும்
முடிவை நிறுவினார். இறுதியாக 1919-இல் கற்பிக்கும் உரிமை அவருக்கு
வழங்கப்பட்டது. பல்கலைக்கழகத் தோழர்களின் தொடர்ந்த எதிர்ப்புக்கு
ஹில்பர்ட், ``கனவான்களே, ஆசிரியர் குழு மன்றம் குளியலறை அல்ல'' என்று
பதிலளித்ததாகக் கூறப்படுகிறது.

% SEGMENT OLP-0626-S03
1920-களின் பிற்பகுதியில் நேதர் அருவ இயற்கணிதத்தில் கவனம் செலுத்தினார்;
அவருடைய பங்களிப்புகள் அத்துறையைப் புரட்சிகரமாக மாற்றின. தமது
நிறுவல்களில் ஏறுவரிசைச் சங்கிலி நிபந்தனை எனப்படுவதை அடிக்கடி
பயன்படுத்தினார்; குறிப்பிட்ட கணங்களின் முடிவிலா, கண்டிப்பாக ஏறும்
சங்கிலி இல்லை என்பதே அதன் பொருள். எடுத்துக்காட்டாக, இப்போது
நேதரியன் வளையங்கள் எனப்படும் சில இயற்கணித அமைப்புகளில்,
இயல்களின் முடிவிலாத் தொடர் $I_1
\subsetneq I_2 \subsetneq \dots$ இல்லை. இந்த நிபந்தனையை எந்தப்
பகுதி வரிசைக்கும் பொதுமைப்படுத்தலாம்; இயற்கணிதத்தில், உட்கணத்
தொடர்பால் வரிசைப்படுத்தப்பட்ட இயல்களே சிறப்பு நிகழ்வு. மறுதலையான
இறங்குவரிசைச் சங்கிலி நிபந்தனையையும் கருதலாம்; அதில் பகுதி வரிசையில்
கண்டிப்பாக \emph{இறங்கு}ம் ஒவ்வொரு தொடரும் முடிவடைகிறது. ஒரு பகுதி
வரிசை இறங்குவரிசைச் சங்கிலி நிபந்தனையை நிறைவு செய்தால்,~$\Nat$ மீதான~$<$
வரிசையில் தொகுத்தறிதலைப் பயன்படுத்துவது போலவே, அந்த வரிசையிலும்
தொகுத்தறிதலைப் பயன்படுத்தலாம். அத்தகைய வரிசைகள்
\emph{நன்கு நிறுவப்பட்டவை} அல்லது \emph{நேதரியன்} எனப்படும்;
அவற்றுக்குரிய நிறுவல் கோட்பாடு \emph{நேதரியன் தொகுத்தறிதல்} எனப்படும்.

% SEGMENT OLP-0626-S04
நேதர் யூதர். 1933-இல் நாசிகள் ஆட்சிக்கு வந்தபோது, அவருடைய பதவியிலிருந்து
நீக்கப்பட்டார். நல்வாய்ப்பாக, பென்சில்வேனியாவின் பிரின் மார் கல்லூரியில்
தற்காலிகப் பணிக்காக அமெரிக்க ஐக்கிய நாடுகளுக்குக் குடிபெயர முடிந்தது.
அங்கிருந்த காலத்தில் பிரின்ஸ்டனிலும் உரையாற்றினார்; ஆனால் பெண்களை
வரவேற்காத பல்கலைக்கழகமாக அதை அவர் உணர்ந்தார் \citep[81]{Dick1981}.
1935-இல் கருப்பையில் இருந்த கட்டியை அகற்ற அறுவை சிகிச்சை செய்துகொண்டார்.
அறுவை சிகிச்சையால் ஏற்பட்ட தொற்றால் இறந்தார்; பிரின் மாரில் அடக்கம்
செய்யப்பட்டார்.

% SEGMENT OLP-0626-S05
\begin{reading}
நேதரின் வாழ்க்கை வரலாற்றுக்கு \citet{Dick1981}-ஐப் பார்க்கவும்.
கோட்பாட்டு இயற்பியலுக்கான பெரிமீட்டர் நிறுவனம், நேதரின் வாழ்க்கை,
தாக்கம் பற்றிய உரைகளை இணையத்தில் வழங்குகிறது \citep{Perimeter2015}.
வாசிப்பிலிருந்து ஓய்வு விரும்பினால், \emph{Stuff You Missed in History
Class} நிகழ்ச்சியில் நேதரின் வாழ்க்கை, தாக்கம் பற்றிய பாட்காஸ்ட்
உள்ளது \citep{Frey2015}. நேதரின் தொகுக்கப்பட்ட படைப்புகள் மூல ஜெர்மன்
மொழியில் கிடைக்கின்றன \citep{Noether1983}.
\end{reading}

\end{document}
